\question[3]
Gegeben ist die folgende Funktion:
\begin{lstlisting}
def rek(n):
    if n == 0: return 0
    return n + rek(n-2)
\end{lstlisting}
\begin{parts}
\part Was liefert \texttt{rek(6)}?
\part Was passiert beim Aufruf \texttt{rek(5)}? Begründe.
\part Ändere den Basisfall so, dass auch \texttt{rek(5)} ein Ergebnis liefert. Welches?
\end{parts}
\begin{solutionbox}{4cm}
\begin{enumerate}[label=\alph*)]
\item $6 + 4 + 2 + 0 = 12$
\item Die Aufrufe sind \texttt{rek(5)}, \texttt{rek(3)}, \texttt{rek(1)}, \texttt{rek(-1)}, \ldots{} Der Basisfall
\texttt{n == 0} wird nie erreicht. Nach zu vielen geschachtelten Aufrufen (normalerweise 1000) bricht Python mit einem \texttt{RecursionError} ab.
\item \texttt{if n <= 0: return 0}. Dann ist \texttt{rek(5)} $= 5 + 3 + 1 + 0 = 9$.
\end{enumerate}
\end{solutionbox}
